\begin{figure}[tb]
\centering
\begin{tikzpicture}[x=0.84cm,y=0.75cm,>=Stealth,
  every node/.style={font=\small}]
 \fill[orange!25] (0,0) rectangle (1,1);
 \fill[blue!9] (1,0) rectangle (10,1);
 \fill[orange!25] (10,0) rectangle (11,1);
 \foreach \x in {0,...,11} {\draw[black!55] (\x,0)--(\x,1);}
 \draw[black!55] (0,0) rectangle (11,1);
 \node[anchor=south,font=\footnotesize] at (0.5,1.08) {nonflat};
 \node[anchor=south,font=\footnotesize] at (10.5,1.08) {nonflat};
 \node[fill=white,inner sep=2pt] at (5.5,0.5)
   {consecutive flat source strips};
 \draw[blue!65!black,thick]
   (1.1,2.75) .. controls (2.1,1.55) and (3.4,1.25) .. (5.5,1.22)
   .. controls (7.6,1.25) and (8.9,1.55) .. (9.9,2.75);
 \node[fill=white,inner sep=2pt] at (5.5,2.35)
   {schematic bound for $m_j=\sup_{S_j}u$};
 \draw[->,blue!65!black] (2.4,1.85)--(3.5,1.42);
 \draw[->,blue!65!black] (8.6,1.85)--(7.5,1.42);
 \draw[<->] (1,-0.35)--(10,-0.35);
 \node[anchor=north] at (5.5,-0.48) {chain length may be arbitrarily large};
 \draw[->] (-0.15,-1.2)--(11.3,-1.2)
   node[anchor=west] {$t$};
 \node[anchor=north,font=\footnotesize] at (5.5,-1.28)
   {source coordinate $v=t+i\theta$};
\end{tikzpicture}
\caption{A flat chain in the source cylinder. The absence of a constant
angular mode makes derivative bounds decay away from the chain's ends.
Their sum, and hence the variation of target height along the chain,
is bounded independently of its length. A chain meeting a target height
slab can therefore be charged to a neighboring nonflat strip in a
fixed enlarged slab. The horizontal coordinate is the source
coordinate $t$, not the target height $h\circ f$.}
\label{fig:flat-chain}
\end{figure}
